\chapter*{คำนำ}
\addcontentsline{toc}{chapter}{คำนำ}

ในยุคดิจิทัลที่เทคโนโลยีสารสนเทศก้าวล้ำ การตรวจติดตามและส่งเสริมสุขภาวะส่วนบุคคล (Personal Health Informatics) มีบทบาทสำคัญอย่างยิ่งต่อการป้องกันโรคไม่ติดต่อเรื้อรัง (Non-Communicable Diseases หรือ NCDs) เช่น โรคเบาหวาน โรคความดันโลหิตสูง และโรคหลอดเลือดหัวใจ ซึ่งเกิดจากพฤติกรรมเนือยนิ่ง (Sedentary Behaviour) เป็นระยะเวลายาวนาน แม้ว่าสมาร์ทโฟนยุคปัจจุบันจะได้รับการติดตั้งตัวตรวจวัดความเร่งแบบสามแกน (3-Axis MEMS Accelerometer) มาตั้งแต่โรงงาน แต่แอปพลิเคชันเพื่อสุขภาพส่วนใหญ่ในท้องตลาดยังคงใช้เกณฑ์การนับก้าวแบบค่าคงที่เชิงเดี่ยว (Static Threshold) ซึ่งมักเกิดความคลาดเคลื่อนสะสมจากการสั่นสะเทือนของยานพาหนะ หรือการสะบัดมือโดยมิได้มีการก้าวเดินจริง

หน่วยวิเคราะห์-ติดตามคุณภาพดินและสภาพอากาศเชิงพื้นที่ คณะวิทยาศาสตร์และเทคโนโลยี มหาวิทยาลัยราชภัฏรำไพพรรณี จึงได้ริเริ่มบูรณาการองค์ความรู้ทางฟิสิกส์จลนศาสตร์ (Kinematics) การประมวลผลสัญญาณดิจิทัล (Digital Signal Processing) เข้ากับหลักวิทยาศาสตร์การกีฬาและชีวกลศาสตร์ (Exercise Science and Biomechanics) เพื่อออกแบบและพัฒนาแอปพลิเคชันมือถือ \textbf{Smart Health Tracker} บนระบบปฏิบัติการข้ามแพลตฟอร์ม Flutter

นวัตกรรมสำคัญของแอปพลิเคชันนี้ประกอบด้วย 3 มิติหลัก ได้แก่
\begin{enumerate}[label=\textbf{\arabic*.}]
  \item \textbf{มิติฟิสิกส์และวิศวกรรมสัญญาณ} การสกัดขนาดเวกเตอร์ความเร่งแบบยูคลิดสามมิติ ($a_{\text{norm}}$) ร่วมกับตัวกรองพีคแบบพลวัต (Dynamic Peak Detection) ที่ปรับระดับเกณฑ์ตามแอมพลิจูดของการเคลื่อนไหวจริง พร้อมทั้งกลไกระยะพักฟื้นทางชีวภาพ (Refractory Period 250 มิลลิวินาที) เพื่อขจัดสัญญาณรบกวน (Artifact Filtering) ได้อย่างแม่นยำ
  \item \textbf{มิติด้านวิทยาศาสตร์สุขภาพเชิงลึก} การคำนวณอัตราความถี่ก้าว (Cadence ในหน่วยก้าวต่อนาที SPM) เพื่อจำแนกความหนักของการออกกำลังกาย (Aerobic Intensity) การคำนวณระยะก้าวแบบแปรผันตามเพศและความสูงของผู้ใช้ การประเมินพลังงานเผาผลาญพื้นฐานและกิจกรรมตามสมการมาตรฐาน Mifflin-St Jeor ตลอดจนระบบแจ้งเตือนการนั่งนานและการดื่มน้ำบริสุทธิ์เพื่อรักษาภาวะสมดุลของเหลวในร่างกาย
  \item \textbf{มิติสถาปัตยกรรมส่วนติดต่อผู้ใช้ยุคใหม่} การจัดวางผังหน้าจอแบบเบนโตะกริด (Bento Grid Layout) ร่วมกับระบบแบ่งส่วนแสดงผลแบบตอบสนองอัตโนมัติตามขนาดจอภาพ (Adaptive Responsive Breakpoints) ที่รองรับทั้งสมาร์ทโฟนหน้าจอแคบ จอมาตรฐาน แท็บเล็ต และคอมพิวเตอร์ตั้งโต๊ะ พร้อมภาพกราฟิกลายเส้นนีออนสไตล์ Cybernetic Health ที่สื่อถึงการผสานเทคโนโลยีเข้ากับสรีรวิทยาได้อย่างลงตัว
\end{enumerate}

คณะผู้จัดทำหวังเป็นอย่างยิ่งว่า คู่มือวิชาการและวิจัยฉบับสมบูรณ์นี้ จะเป็นประโยชน์ต่อคณาจารย์ นักศึกษา นักพัฒนาซอฟต์แวร์ และผู้สนใจทั่วไป ในการนำไปประยุกต์ใช้ในการเรียนการสอน การวิจัย และการจัดกิจกรรมอบรมเชิงปฏิบัติการ เพื่อส่งเสริมสุขภาวะที่ดีของสังคมไทยอย่างยั่งยืนสืบไป

\vspace{1.5cm}

\begin{flushright}
  \textbf{ผู้ช่วยศาสตราจารย์ ดร.ชีวะ ทัศนา}\\
  หน่วยวิเคราะห์-ติดตามคุณภาพดินและสภาพอากาศเชิงพื้นที่\\
  คณะวิทยาศาสตร์และเทคโนโลยี มหาวิทยาลัยราชภัฏรำไพพรรณี
\end{flushright}
